\subsection{Topological Spaces, Bases, and Separation}

% **Role:** abstract structure behind limits and continuity.
%
% Topological spaces, bases, product topology, quotient topology, countability axioms.


General topology abstracts the open-set behavior of metric spaces. This selective pass records spaces, bases, products, quotients, and countability axioms for later use.

\begin{definition} \label{def:topological-space} \hlt{Topological Space}\\
A \hlt{topology} on a set $X$ is a collection $\mathcal{T}\subseteq 2^X$ such that
\begin{enumerate}[label=(\roman*)]
\setlength{\itemsep}{0pt}
\item $\emptyset,X\in\mathcal{T}$;
\item arbitrary unions of members of $\mathcal{T}$ belong to $\mathcal{T}$;
\item finite intersections of members of $\mathcal{T}$ belong to $\mathcal{T}$.
\end{enumerate}
The pair $(X,\mathcal{T})$ is a \hlt{topological space}, and the members of $\mathcal{T}$ are the \hlt{open sets}. Closed sets are complements of open sets.
\end{definition}

\begin{proposition} \label{prop:metric-open-sets-form-topology} \hlt{Metric Open Sets Form a Topology}\\
If $(X,d)$ is a metric space, then the open sets from Definition \ref{def:metric-topology-core-definitions} form a topology on $X$.
\begin{proof}
The empty set and $X$ are open. Arbitrary unions and finite intersections are open by Theorem \ref{thm:metric-open-closed-set-calculus}.\\
\hlt{Structural Commentary}: Metric balls generate an open-set system with the exact axioms of a topology.\\
\hlt{Key Insight}: General topology keeps the open-set calculus and forgets the distance function.
\end{proof}
\end{proposition}

\begin{definition} \label{def:basis-for-topology} \hlt{Basis}\\
Let $X$ be a set. A collection $\mathcal{B}\subseteq 2^X$ is a \hlt{basis} for a topology if:
\begin{enumerate}[label=(\roman*)]
\setlength{\itemsep}{0pt}
\item for every $x\in X$, some $B\in\mathcal{B}$ contains $x$;
\item if $x\in B_1\cap B_2$ with $B_1,B_2\in\mathcal{B}$, then some $B_3\in\mathcal{B}$ satisfies $x\in B_3\subseteq B_1\cap B_2$.
\end{enumerate}
The topology generated by $\mathcal{B}$ consists of all unions of basis elements. A \hlt{subbasis} is a collection whose finite intersections form a basis.
\end{definition}

\begin{example} \label{ex:metric-balls-basis} \hlt{Metric Balls as a Basis}\\
In a metric space, the collection of all open balls is a basis for the metric topology. On $\R^k$, balls with rational centers and rational radii form a countable basis.
\end{example}
